%=============================================================================
%  A1_reproducibility.tex -- Appendix A: reproducibility
%
%  CONVENTION: hardware identifiers, driver and library version strings, and
%  file paths are provenance identifiers, not measured quantities, and are
%  written literally. Every measured quantity elsewhere in the paper is a macro
%  from macros.tex.
%=============================================================================
\section{Reproducibility}
\label{sec:repro}

\subsection{Experimental campaigns}

The results in this paper come from four campaigns run on \emph{three different
hosts with three different software stacks}. They are not interchangeable, and
no table in this paper mixes them. \Cref{tab:campaigns} is the authority for
which campaign produced which result; every results table elsewhere names its
campaign in the caption.

\begin{table}[t]
  \centering
  \small
  \caption{The four experimental campaigns. Data throughout is the BGL corpus
  from LOGHUB, temporal split, test partition, BF16 precision.}
  \label{tab:campaigns}
  \begin{tabular}{llp{0.20\linewidth}p{0.24\linewidth}p{0.20\linewidth}}
    \toprule
    Campaign & Issues & GPU / driver & Runtime & Models\\
    \midrule
    A --- AR baseline correction
      & \#12
      & H100 PCIe 80GB, driver 580.65.06 (CUDA 13.0)
      & vLLM 0.19.1, torch 2.10.0+cu128
      & Qwen2.5-1.5B / 3B / 7B\\[2pt]
    B --- cost laws, walls, applications
      & \#14, \#15, \#17, \#19, \#21
      & A100 PCIe 80GB, CUDA-13 driver
      & vLLM 0.25.1, torch 2.11.0+cu130
      & \modelDG, \modelARthree\\[2pt]
    C --- diffusion-native serving control
      & \#22, \#23
      & A100 PCIe 80GB
      & dInfer 0.1 + vLLM 0.10.2, torch 2.8.0+cu128, transformers 4.57.6
      & \modelLLaDA\\[2pt]
    D --- original HF-transformers pass
      & \#3--\#8
      & H100 PCIe 80GB
      & HF transformers 5.14.1, torch 2.5.1+cu121
      & Qwen2.5-1.5B / 7B, LLaDA-8B\\
    \bottomrule
  \end{tabular}
\end{table}

Three consequences of this split are load-bearing. First, the corrected AR
baseline (Campaign~A) and the original pass it corrects (Campaign~D) ran on the
same GPU model and the same 2000 BGL events, so the before-and-after comparison
is a comparison of stacks and decoding settings, not of hardware. Second, the
cost laws, the two walls, the batching multipliers and all of
\cref{sec:applications} come from a single campaign (B) on a single host, so
they are internally comparable even though they are not comparable to
Campaign~A's absolute latencies. Third, the dInfer control that licenses the
architectural reading of the batching cap (\cref{sec:mechanism}) is Campaign~C,
a different runtime on a different stack --- which is precisely why it is
evidence, and also why its latencies must not be compared against Campaign~B's.

\subsection{Claim-to-artifact map}

Every numeric claim in this paper is traceable to a committed JSON artifact
under \texttt{results/}. \Cref{tab:artifacts} maps sections to artifacts. The
values themselves are collected in \texttt{paper/NUMBERS.md}, which was
assembled by reading these JSON files directly rather than the intermediate
Markdown reports; where a report disagreed with its own JSON, the JSON is
authoritative and the disagreement is recorded in that file. The script
\texttt{paper/verify\_numbers.py} re-reads every artifact and asserts each
recorded value; it is intended to run before every submission build.

\begin{table}[t]
  \centering
  \small
  \caption{Claim-to-artifact map. All paths are relative to the repository root.}
  \label{tab:artifacts}
  \begin{tabular}{llp{0.46\linewidth}}
    \toprule
    Claim & Campaign & Artifact\\
    \midrule
    AR baseline correction (\cref{sec:method})
      & A & \texttt{results/vllm\_serving\_\{1.5b,3b,7b\}.json}\\
    Fitted cost laws (\cref{sec:cost-laws})
      & B & \texttt{results/appD\_analysis.json}\\
    Quality invariance to output length
      & B & \texttt{results/p2\_\{dg,ar3b\}\_L\{32,64,128,256,512\}.json},
            \texttt{results/appD\_\{dg,ar3b\}\_L\{1024,2048\}.json}\\
    Two orthogonal walls (\cref{sec:two-walls})
      & B & \texttt{results/rate\_lag\_\{dg,ar3b\}*.json};
            engine crash: \texttt{results/rate\_lag\_dg\_engine\_crash.json}\\
    Batching multipliers (\cref{sec:mechanism})
      & B, C & \texttt{results/appD\_mult\_\{dg,ar3b\}\_L\{128,512\}.json},
            \texttt{results/llada\_dinfer/}\\
    Serial chains (\cref{sec:app-chains})
      & B & \texttt{results/chain\_\{dg,ar3b\}.json},
            \texttt{results/chain\_analysis.json}\\
    Wide schema (\cref{sec:app-wide})
      & B & \texttt{results/wide\_schema\_\{dg,ar3b\}.json},
            \texttt{results/wide\_schema\_analysis.json}\\
    Bounded deadline (\cref{sec:app-deadline})
      & B, D & \texttt{results/bounded\_deadline\_\{ar3b,dm\_hf\}.json}\\
    Iterative refinement (\cref{sec:app-refine})
      & B & \texttt{results/refinement\_report.md}\\
    Priority preemption (\cref{sec:app-vip})
      & B & \texttt{results/vip\_\{dg,ar3b\}.json}\\
    Perturbation robustness (\cref{sec:neg-perturb})
      & B, D & \texttt{results/perturbation\_\{ar3b,dm\_hf\}.json}\\
    One-shot imputation (\cref{sec:neg-impute})
      & B, D & \texttt{results/imputation\_\{AR,DM\}\_random.json}\\
    Adaptive effort (\cref{sec:neg-adaptive})
      & B, C & \texttt{results/adaptive\_effort\_\{ar3b,dinfer\}.json},
            \texttt{results/difficulty\_bins.json}\\
    Semi-AR blocks, $K$-sweep (\cref{sec:neg-semiar,sec:neg-survived})
      & D & \texttt{results/ablation\_table5.json}\\
    Classical baselines (\cref{sec:neg-baselines})
      & D & \texttt{results/baselines\_scenarios.json},
            \texttt{results/per\_field\_fig4.json}\\
    dInfer control (\cref{sec:threats})
      & C & \texttt{results/llada\_dinfer/}\\
    \bottomrule
  \end{tabular}
\end{table}

\subsection{Runners}

Every experiment in this paper is driven by a committed runner under
\texttt{models/experiments/}, and every analysis by a committed analysis script
in the same directory. The naming is mechanical: \texttt{run\_*.py} produces the
JSON artifacts in \cref{tab:artifacts} and \texttt{analyze\_*.py} reduces them to
the tables reported here --- for example \texttt{run\_chain\_experiment.py} and
\texttt{analyze\_chains.py} for \cref{sec:app-chains}, or
\texttt{run\_perturbation.py} and \texttt{analyze\_perturbation.py} for
\cref{sec:neg-perturb}. No number in this paper was produced by an
uncommitted script or by hand.

\subsection{What we did not run}

FlowFill's weights are not released. We never executed it. Every reference to it
in this paper is a reference to its published claims; its absolute figures ---
its reported latency, its reported throughput, and its reported speedup over
autoregressive serving --- are untested here and are excluded from
\texttt{paper/NUMBERS.md} for that reason. We have not reproduced them and we
have not failed to reproduce them. Readers should not infer from any comparison
in this paper that its numbers were checked; what was checked is what our own
substitute models do on our own hardware under our own harness, which is a
weaker and different statement.
